\begin{table}[!htbp]
\centering
\caption{Symmetrized (Shapley/Oaxaca--Blinder) decomposition of mean differences between Control and Market into a representation component (differences in the distribution of social proximity) and a behavior component (differences in allocations conditional on social proximity), dictator games pooled, hypothetical-allocation sample. The decomposition is a descriptive accounting; no sampling uncertainty is attached.}
\label{tab:decomposition_sym_pooled_m0m1}
\begin{tabular}{lccc}
\toprule
\textbf{Variable} & \textbf{Observed Diff} & \textbf{Representation Component} & \textbf{Behavior Component} \\
\midrule
\multicolumn{4}{l}{\textbf{Control $-$ Market}} \\
Mean Allocation    & $-$1.996 & $-$0.494 & $-$1.502 \\
Allocation = 4     & 0.025 & $-$0.002 & 0.027 \\
Allocation = 6     & 0.416 & 0.097 & 0.319 \\
Allocation = 8     & 0.028 & $-$0.012 & 0.040 \\
Allocation $>8$    & $-$0.461 & $-$0.083 & $-$0.378 \\
Allocation = 12    & $-$0.150 & $-$0.047 & $-$0.103 \\
\bottomrule
\end{tabular}
\begin{flushleft}
\footnotesize Notes: The construction is the symmetrized decomposition of Section~\ref{sec:heterogeneity}, with representation cells given by the five social-proximity levels of Table~\ref{tab:freqs_market_control_sp} ($N=400$ Control, $N=800$ Market); no cell is empty in either condition. Cells defined by the seven moral categories instead split the mean-allocation gap into $-$0.617 (representation) and $-$1.379 (behavior).
\end{flushleft}
\end{table}
